\documentclass[10pt,a4paper]{article}

%double si on veut une version prof et une version élève. simple si on veut une seule version.
\newcommand{\buildmode}{double}

\InputIfFileExists{version.tex}{}{}
% Version (définie par le script)
\providecommand{\version}{eleve}

\usepackage{feuilleactivite}

%%%à modifier à chaque fois

\renewcommand{\niveau}{2nde}
\renewcommand{\chapitre}{Equations produit nul}
\renewcommand{\numchap}{0.4} 
\renewcommand{\theme}{Fonctions}

\begin{document}

\prof{
    Motivation : Introduire la fonction carré comme une grandeur non proportionnelle. Si je multiplie le côté d'un carré par 2, est-ce que son aire est multipliée par 2 ?\\
    On peut faire des exemples au tableau puis la démonstration. Ensuite, faire tracer la fonction carré. On peut préparer une activité pour ça.
}
\begin{objectifs}
\begin{itemize}
    \item Connaitre la représentation graphique de la fonction carré
    \item Savoir résoudre des équations et des inéquations avec la fonction carré
\end{itemize}
\end{objectifs}

\begin{multicols}{2}

    \begin{exo}
        On a représenté la fonction carré
    \end{exo}

\end{multicols}
\end{document}
